\documentclass{article}
\usepackage{graphicx} % Required for inserting images
\usepackage{amsmath}
\usepackage{todonotes} % temporario
\usepackage{amssymb} % para usar \mathbb{R}

\begin{document}

\section{Introdução}

\par O estudo de redes constitui um problema clássico em diversas áreas da engenharia, física e matemática aplicada. Neste contexto, considera-se uma rede hidráulica composta por nós e arestas, na qual cada aresta possui uma resistência hidráulica associada e cada nó pode atuar como fonte, sumidouro ou ponto de pressão conhecida.

No problema direto abordado neste trabalho, assume-se que uma vazão $Q_{in}$ é injetada em um ou mais nós da rede, enquanto outros nós estão submetidos a uma condição de contorno de pressão conhecida, tipicamente a pressão atmosférica. O objetivo consiste em determinar as pressões em todos os nós e as vazões em todas as arestas, a partir das leis de conservação de massa e das relações constitutivas que governam o escoamento.

Além das variáveis primárias do problema, também são de interesse diversas grandezas derivadas que caracterizam o comportamento global da rede, tais como a resistência hidráulica equivalente, a potência necessária para sustentar o escoamento imposto e métricas de eficiência associadas à dissipação de energia no sistema.

A resolução eficiente deste problema requer a formulação de um sistema linear que represente o balanço de fluxos nos nós e as relações entre pressões e vazões nas arestas. Neste trabalho, tal sistema é construído utilizando uma abordagem matricial baseada em operadores de incidência e condutância, resultando em um sistema esparso que pode ser resolvido de forma computacionalmente eficiente.

Para fins de análise e visualização, foi desenvolvida uma ferramenta computacional interativa em Python, utilizando as bibliotecas \texttt{NumPy}, \texttt{SciPy} e \texttt{Pygame}. Essa ferramenta permite a geração automática de redes aleatórias, bem como a manipulação direta de nós e arestas, possibilitando a exploração dinâmica do comportamento do sistema.

\section{Desenvolvimento}

\subsection{Problema 1}
\indent\indent Suponha que à rede seja injetada uma vazão $Q_{in}$ em um ou mais pontos, enquanto outros nós estejam submetidos à pressão atmosférica.
O problema consiste em determinar: 
\begin{itemize}
    \item as vazões em cada aresta;
    \item as pressões em cada nó;
    \item a resistência hidráulica equivalente da rede;
    \item a potência necessária para injetar a vazão $Q_{in}$;
    \item outras métricas globais do sistema. 
\end{itemize}

\subsubsection{Formulação Matemática Método Direto}
\indent\indent Dado um grafo $G$, em que $V$ é o conjunto de nós da rede e $A$ é o conjunto de arestas (canos) conectando os nós:

\[
G = (V, A), \quad |V| = n, \quad |A| = m
\]

O grafo $G$ se torna uma rede se estiver munido de duas regras adicionais:

\begin{equation}
\sum_{k \in \delta(j)} q_k = Q_j, \quad j = 1, ..., n
\label{eq:lei_constitutiva}
\end{equation}

(Lei de constituição)

\todo{colocar os nomes das leis ao lado}


\begin{equation}
q_k = \frac{1}{R_k} \, (p_i - p_j), \quad k=(i,j) \in A
\label{eq:lei_balanco}
\end{equation}
(Lei de equilíbrio)

Onde:

\[
 p \in \mathbb{R}^n, \quad
p =
\begin{bmatrix}
p_1 \\
p_2 \\
\vdots \\
p_n
\end{bmatrix}
\]

\[
q \in \mathbb{R}^m, \quad
q =
\begin{bmatrix}
q_1 \\
q_2 \\
\vdots \\
q_m
\end{bmatrix}
\]

\[
Q \in \mathbb{R}^n, \quad
Q =
\begin{bmatrix}
Q_1 \\
Q_2 \\
\vdots \\
Q_n
\end{bmatrix}
\]

\begin{itemize}
    \item $p_i$ é a pressão no vértice $i$, variável dependente;
    
    \item $q_k$ é a vazão na aresta $k$, variável dependente; 

    \item $Q_i$ é a vazão no vértice $i$. e $Q_{in}$ é a vazão que entra no sistema; 
    
    \item $\delta(j) = \{ k \in A \mid k \text{ é incidente ao nó } j \}$;
    
    %\item $a_{ij} = 1$, se vértice $i$ está ligado ao $j$, senão $0$. 

    \item $R_k$ é a resistência ao fluxo, na aresta $k$. $R_k >0$

    %\item $Q_{in}$ é a vazão de entrada
\end{itemize}

Deste modo, as equações constitutivas \eqref{eq:lei_constitutiva} e de balanço nodal \eqref{eq:lei_balanco} formam o seguinte sistema de equações:

\begin{equation}
\left\{
\begin{aligned}
    \sum_{k \in \delta(1)} q_k &= Q_1 \\
    &\vdots \\
    \sum_{k \in \delta(n)} q_k &= Q_n \\
    q_1 - \frac{1}{R_1}(p_{i_1} - p_{j_1}) &= 0 \\
    &\vdots \\
    q_m - \frac{1}{R_m}(p_{i_m} - p_{j_m}) &= 0
\end{aligned}
\right.
\label{eq:sistema}
\end{equation}

Observa-se que o sistema é composto por \(m+n\) equações e \(m+n\) incógnitas, sendo \(q_k\) as vazões nos ramos da rede e \(p_i\) as pressões nos nós. Dessa forma, o sistema \eqref{eq:sistema} pode ser reescrito de forma compacta na forma matricial

\[
Zx = b,
\]

onde o vetor de incógnitas é definido por

\[
x =
\begin{bmatrix}
q_1 \\
\vdots \\
q_m \\
p_1 \\
\vdots \\
p_n
\end{bmatrix},
\quad
x \in \mathbb{R}^{m+n},
\]

e o vetor de termos independentes é dado por

\[
b =
\begin{bmatrix}
Q_1 \\
\vdots \\
Q_n \\
0 \\
\vdots \\
0
\end{bmatrix},
\quad
b \in \mathbb{R}^{m+n}.
\]

A matriz de coeficientes \(Z\) pode ser estruturada em blocos da seguinte forma:

\[
Z =
\begin{bmatrix}
C & \textbf{0} \\
I & D
\end{bmatrix}_ {(m+n)\times (m+n)}.
\]


O bloco \(C \in \mathbb{R}^{n \times m}\) corresponde às equações de balanço nodal e depende apenas da estrutura do grafo que representa a rede. Seus elementos assumem valores \(0\) ou \(1\), indicando se um determinado ramo incide em um dado nó.

O bloco superior direito é uma matriz nula \(0_{n \times n}\), pois as equações de balanço nodal não dependem diretamente das pressões nos nós.

O bloco inferior esquerdo é a matriz identidade \(I_m\), uma vez que cada equação constitutiva possui coeficiente unitário associado à vazão \(q_k\).

Por fim, o bloco \(D \in \mathbb{R}^{m \times n}\) contém os coeficientes associados às pressões nas equações constitutivas. Esses coeficientes dependem tanto da estrutura do grafo quanto das resistências hidráulicas \(R_k\) de cada ramo, sendo compostos por termos da forma \(\pm \frac{1}{R_k}\) associados aos nós conectados pelo ramo \(k\).

Assim, o sistema completo pode ser representado como

\[
\begin{bmatrix}
C & \textbf{0} \\
I & D
\end{bmatrix}
\begin{bmatrix}
q \\
p
\end{bmatrix}
=
\begin{bmatrix}
Q \\
0
\end{bmatrix}.
\]

Além disso, para o problema admitir solução única, é necessário que $\det Z \neq 0$, pois senão o sistema poderá ser singular, o que não é desejado. 

Para garantir a não-singularidade, é necessário fixar uma pressão atmosférica em algum vértice, garantindo uma solução para o problema, o que é feito modificando os blocos de matriz $C$ e $\textbf{0}$, superiores à matriz $Z$. 

Para isso, se fixada pressão atmosférica ("cano furado") no vértice $i$ então toda a linha $i$ de $C$ será definida como 0, obtemos assim a matriz $C_{mod}$. Além disso a matriz $\textbf{0}$ recebe 1 na posição diagonal $(i,i)$, tornando-se $\textbf{0}_{mod}$. 

Então em sua forma final:


\[
\begin{bmatrix}
C_{mod} & \textbf{0}_{mod} \\
I & D
\end{bmatrix}
\begin{bmatrix}
q \\
p
\end{bmatrix}
=
\begin{bmatrix}
Q \\
0
\end{bmatrix}.
\]



% A partir desse sistema, é possível reescrevê-lo como matriz da seguinte forma:

% Definindo o vetor de incógnitas

% \[
% x =
% \begin{bmatrix}
% q_1 \\
% \vdots \\
% q_m \\
% p_1 \\
% \vdots \\
% p_n
% \end{bmatrix},
% x \in \mathbb{R}^{m+n}
% \]

% e o vetor de resultado
% \[
% b =
% \begin{bmatrix}
% Q_1 \\
% \vdots \\
% Q_n \\
% 0 \\
% \vdots \\
% 0
% \end{bmatrix},
% b \in \mathbb{R}^{m+n}
% \]



% tem-se

%\[
% \begin{bmatrix}
% 1 & 0 & 0 \\
% -1 & 0 & 0 \\
% 1 & -\frac{1}{R_1} & \frac{1}{R_1}
% \end{bmatrix}
% \begin{bmatrix}
% q_1 \\
% p_1 \\
% p_2
% \end{bmatrix}
% =
% \begin{bmatrix}
% Q_1 \\
% Q_2 \\
% 0
% \end{bmatrix}.
% \]


\paragraph{Resistência Hidráulica Equivalente da Rede}
A resistência hidráulica equivalente representa o quanto a rede dificulta o escoamento do fluido. Esse valor pode ser obtido a partir da razão entre a diferença de pressão nos extremos da rede e a vazão de entrada no sistema.

\begin{equation}
R_{eq} = \frac{p_1 - p_n}{Q_{in}}
\label{eq:Req}
\end{equation}

onde $Q_{in}$ é a vazão de entrada, $p_n$ é a pressão no último ponto da rede, e $p_1$ é a pressão no ponto de entrada. Fisicamente, $R_{eq}$ corresponde à resistência de um único tubo equivalente que produziria a mesma queda de pressão para a vazão $Q_{in}$ que toda a rede original.

\paragraph{Potência necessária para injetar a vazão de entrada} 
A potência hidráulica necessária para injetar a vazão $Q_{in}$ ao longo de toda a rede é equivalente à energia total dissipada nos tubos do sistema, associada às perdas de pressão em cada trecho da tubulação. Essa potência pode ser expressa por:

\begin{equation}
    P_{\text{in}} = \sum_{k} q_k \cdot \Delta p_k
\end{equation}

onde $k$ representa uma aresta, $q_k$ é a vazão que escoa pela aresta $k$, e $\Delta p_k$ é a diferença de pressão entre as extremidades dessa aresta. 

\par O método é dito direto pois é montada apenas um sistema de matrizes para encontrar a solução

\paragraph{Pré Processamento e Assimilação de Dados}


Após a formulação matemática, é necessário instanciar
valores concretos para os parâmetros da rede e preparar os dados que
serão utilizados na resolução computacional.

Inicialmente, define-se o formato da rede a partir do grafo $G$, especificando quais pares
de vértices estão conectados por arestas. 

Para cada aresta $k \in A$ atribui-se um valor de resistência $R_k$.

Também são definidos os valores das vazões nodais $Q_i$.

Nós que recebem injeção de fluido possuem $Q_i>0$, enquanto nós de
descarga possuem $Q_i<0$. Nós internos da rede satisfazem $Q_i=0$, pois tudo que entra, sai.
A vazão total injetada no sistema é denotada por $Q_{in}$.

Utilizaremos a pressão atmosférica como referência para o sistema,
fixando $p_{atm} = 0$.

Com essas informações, é possível construir o sistema linear que
determina as pressões da rede. Em sua forma simples:

\[
Zx = b ,
\]

no qual a matriz $Z$ depende apenas da topologia da rede e dos valores
das resistências $R_k$. A partir da solução desse sistema, obtêm-se as
pressões nodais e, posteriormente, as vazões em cada aresta por meio
da lei constitutiva.

Após a construção do sistema, é necessário verificar se o problema é bem-posto (Hadamard), ou seja, se este tem solução e essa solução é única. Para isso, é necessário que ocorra  $det(Z)\neq0$.


\subsubsection{Formulação Matemática Método de Duas Etapas}

Uma abordagem alternativa à formulação matricial acoplada consiste no método do qual as \textbf{vazões nos ramos são determinadas
antes do cálculo das pressões nodais}. Essa técnica baseia-se na aplicação simultânea da conservação de massa nos nós e da conservação de energia nos ciclos independentes do grafo.

Inicialmente, escolhe-se um conjunto de vazões iniciais $q \in \mathbb{R}^m$
que satisfaça a equação de continuidade nodal

\begin{equation}
C\,q = Q,
\end{equation}

onde $C \in \mathbb{R}^{n \times m}$ é a matriz de incidência nó--aresta
e $Q \in \mathbb{R}^n$ representa as vazões impostas nos nós.

Em seguida, em vez de escrever equações nos nós, escrevem-se equações
nos \textbf{laços independentes} da rede. Seja $B \in \mathbb{R}^{\ell \times m}$
a \textbf{matriz de laços} (ou matriz ciclo--aresta), onde $\ell = m-n+1$
é o número de ciclos independentes do grafo. Cada linha de $B$ indica,
com sinais $\pm 1$, quais arestas pertencem a um determinado laço e
o respectivo sentido de percurso.

A conservação de energia em cada laço impõe que a soma das perdas de
pressão ao longo do ciclo seja nula, resultando em

\begin{equation}
B\,h = 0,
\end{equation}

onde $h \in \mathbb{R}^m$ é o vetor de perdas de carga nas arestas,
definido por

\begin{equation}
h = R\,q,
\end{equation}

sendo $R \in \mathbb{R}^{m \times m}$ a matriz diagonal das resistências
hidráulicas,

\[
R =
\begin{bmatrix}
R_1 & 0 & \cdots & 0 \\
0 & R_2 & \cdots & 0 \\
\vdots & \vdots & \ddots & \vdots \\
0 & 0 & \cdots & R_m
\end{bmatrix}.
\]

Assim, as equações dos laços podem ser escritas de forma compacta como

\begin{equation}
B\,R\,q = 0.
\end{equation}

O método procede então de forma iterativa, aplicando correções às vazões
de modo a satisfazer essas equações, sem a necessidade de montar um
grande sistema linear acoplado.

Após a convergência das vazões, as pressões nodais são obtidas por
integração direta das quedas de pressão ao longo das arestas. Fixando-se
a pressão em um nó de referência $p_{\text{ref}} = 0$, as demais pressões
podem ser calculadas por

\begin{equation}
p_j = p_i - R_k\,q_k,
\quad \text{para cada aresta } k = (i,j).
\end{equation}

Dessa forma, diferentemente da formulação matricial global, o método
dos laços resolve inicialmente apenas as vazões, utilizando um número de
equações igual ao número de ciclos independentes da rede,
$\ell = m-n+1$, que é significativamente menor que $m+n$ em redes
de grande porte. Isso torna essa abordagem particularmente eficiente
para redes extensas e esparsas.


\subsection{Problema 2}

Considere a mesma rede hidráulica do Problema Determinístico,
representada por um grafo \( G=(V,A) \), com \( |V|=n \) nós e \( |A|=m \) arestas.
No caso determinístico, cada aresta \(k\) possui resistência hidráulica fixa \(R_k\).
Neste novo cenário, admite-se que cada cano da rede pode estar parcialmente obstruído
com probabilidade \(r\), o que multiplica sua resistência por um fator \(\alpha>1\).

Para modelar essa incerteza, define-se, para cada aresta \(k\), a variável aleatória discreta
\[
\xi_k =
\begin{cases}
1, & \text{com probabilidade } 1-r, \\[6pt]
\alpha, & \text{com probabilidade } r.
\end{cases}
\]

A resistência hidráulica passa a ser uma variável aleatória
\[
\widetilde{R}_k = \xi_k R_k.
\]

Logo, a matriz diagonal das resistências da rede torna-se
\[
\widetilde{R} =
\begin{bmatrix}
\widetilde{R}_1 & 0 & \cdots & 0 \\
0 & \widetilde{R}_2 & \cdots & 0 \\
\vdots & \vdots & \ddots & \vdots \\
0 & 0 & \cdots & \widetilde{R}_m
\end{bmatrix}.
\]

\subsubsection{Formulação Matemática}

Seja \(C \in \mathbb{R}^{n \times m}\) a matriz de incidência nó–aresta da rede.
A relação entre vazões nas arestas \(q \in \mathbb{R}^m\) e pressões nodais
\(p \in \mathbb{R}^n\) é dada por
\[
q = D(\widetilde{R})\, C^T p,
\]
onde
\[
D(\widetilde{R}) =
\begin{bmatrix}
\frac{1}{\widetilde{R}_1} & 0 & \cdots & 0 \\
0 & \frac{1}{\widetilde{R}_2} & \cdots & 0 \\
\vdots & \vdots & \ddots & \vdots \\
0 & 0 & \cdots & \frac{1}{\widetilde{R}_m}
\end{bmatrix}.
\]

A conservação de massa nos nós impõe
\[
C q = Q,
\]
onde \(Q \in \mathbb{R}^n\) é o vetor de vazões injetadas/retiradas.

Substituindo a expressão de \(q\), obtém-se o sistema reduzido
em termos apenas das pressões nodais:
\[
C\, D(\widetilde{R})\, C^T\, p = Q.
\]

Define-se então a matriz Laplaciana hidráulica aleatória
\[
L(\widetilde{R}) = C\, D(\widetilde{R})\, C^T.
\]

O sistema a ser resolvido passa a ser
\[
L(\widetilde{R})\, p = Q.
\]

Como \(\widetilde{R}\) depende das variáveis aleatórias \(\xi_k\),
a matriz \(L\) é aleatória e, consequentemente,
o vetor de pressões
\[
p = p(\xi_1,\xi_2,\dots,\xi_m)
\]
também é uma variável aleatória vetorial.

O objetivo do problema é determinar
\[
\mathbb{P}\!\left(\max_{i=1,\dots,n} p_i > P_{\max}\right),
\]
isto é, a probabilidade de que a pressão em algum nó ultrapasse
o valor máximo admissível.

\subsubsection{Método de Monte Carlo}

Como cada aresta possui dois estados possíveis,
o número total de configurações da rede é \(2^m\),
o que inviabiliza uma análise exata para redes reais.

Utiliza-se então o Método de Monte Carlo, baseado em amostragem aleatória:

\begin{enumerate}
\item Para \(s = 1, \dots, N\), gera-se uma amostra independente
\(\{\xi_k^{(s)}\}_{k=1}^m\).

\item Para cada amostra, monta-se a matriz
\(\widetilde{R}^{(s)}\), a matriz
\(D(\widetilde{R}^{(s)})\) e a Laplaciana
\[
L^{(s)} = C\, D(\widetilde{R}^{(s)})\, C^T.
\]

\item Resolve-se o sistema linear
\[
L^{(s)} p^{(s)} = Q.
\]

\item Verifica-se a condição de falha
\[
I^{(s)} =
\begin{cases}
1, & \text{se } \max_i p_i^{(s)} > P_{\max}, \\
0, & \text{caso contrário.}
\end{cases}
\]
\end{enumerate}

A probabilidade procurada é estimada por
\[
\mathbb{P}\!\left(\max_i p_i > P_{\max}\right)
\approx
\frac{1}{N} \sum_{s=1}^{N} I^{(s)}.
\]

O problema transforma a análise hidráulica da rede em um problema
de propagação de incertezas em sistemas lineares dependentes de parâmetros
aleatórios. A obstrução de canos altera as resistências, modifica a
distribuição de vazões e pode provocar acúmulo de pressão em determinados nós.

A formulação reduzida em pressões nodais é essencial para tornar viável
a simulação repetida do sistema, já que o método exige a resolução de
centenas ou milhares de sistemas lineares.

Essa abordagem caracteriza um estudo de confiabilidade da rede hidráulica,
permitindo quantificar o risco de violação dos limites operacionais
de pressão devido a falhas aleatórias nas arestas.




% \subsubsection{Métodos e Implementação Computacional}

% linguagem: python


% montagem matriz


% método de solução: cholesky, L*U, etc



% --

% \subsubsection{Pós-Processamento e Análise de Resultados}


% \subsubsection{Variáveis Dependentes}

% \begin{itemize}
%     \item Vazão nas arestas
%     \item Pressão em cada nó
%     \item Potência hidráulica ($R_{eq}$): $R_{eq} = \frac{P_0 - P_{in}}{Q_{in}}$, informação da rede
%     \item Diâmetro do cano
%     \item Densidade/viscosidade do líquido
%     \item Temperatura do líquido
%     \item $R_{eq}$
% \end{itemize}
 
% \subsubsection{Condições de Contorno}

% \begin{itemize}
%     \item Pressão atmosférica fixa $p_a$
% \end{itemize}


\section{Método de Monte Carlo}

O método de Monte Carlo é uma técnica numérica baseada em amostragem aleatória para estimar quantidades de interesse em problemas determinísticos ou estocásticos. A ideia central consiste em modelar as incertezas do sistema por meio de variáveis aleatórias, gerar diversas realizações independentes dessas variáveis e, para cada realização, calcular a grandeza de interesse. A média dos resultados obtidos converge para o valor esperado da variável analisada.

No presente problema, cada aresta da rede pode sofrer um entupimento com probabilidade \(p\). Quando isso ocorre, a resistência hidráulica da aresta é multiplicada por um fator \(\alpha > 1\). Assim, para cada simulação, gera-se um vetor aleatório de resistências
\[
R^{(k)} = \left(R_1^{(k)}, R_2^{(k)}, \dots, R_m^{(k)}\right),
\]
onde
\[
R_i^{(k)} =
\begin{cases}
\alpha R_i, & \text{com probabilidade } p,\\
R_i, & \text{com probabilidade } 1-p.
\end{cases}
\]

Para cada realização \(R^{(k)}\), resolve-se o sistema linear da rede e obtêm-se as pressões nos nós,
\[
p^{(k)} = \left(p_1^{(k)}, p_2^{(k)}, \dots, p_n^{(k)}\right).
\]
Em seguida, calcula-se a pressão máxima
\[
P_{\max}^{(k)} = \max_i p_i^{(k)}.
\]

Sejam \(N\) simulações independentes. A média amostral
\[
\overline{P}_{\max}
=
\frac{1}{N}
\sum_{k=1}^{N}
P_{\max}^{(k)}
\]
é um estimador do valor esperado \(\mathbb{E}[P_{\max}]\).

De modo análogo, a probabilidade de a pressão máxima ultrapassar um limite \(L\) pode ser estimada por
\[
\widehat{P}
=
\frac{1}{N}
\sum_{k=1}^{N}
\mathbf{1}\!\left(P_{\max}^{(k)} > L\right),
\]
onde \(\mathbf{1}(\cdot)\) é a função indicadora.

\section{Erro do Método de Monte Carlo}

O método de Monte Carlo possui convergência assintótica governada pela Lei dos Grandes Números e pelo Teorema Central do Limite. O erro estatístico da média amostral é proporcional a
\[
\mathcal{O}\!\left(\frac{1}{\sqrt{N}}\right).
\]

Mais precisamente, se \(\sigma^2\) é a variância da variável aleatória de interesse, então o erro padrão da média é
\[
\text{EP}
=
\frac{\sigma}{\sqrt{N}}.
\]

No caso da estimativa de uma probabilidade, o erro padrão é dado por
\[
\text{EP}
=
\sqrt{
\frac{\widehat{P}(1-\widehat{P})}{N}
}.
\]

Um intervalo de confiança aproximado de \(95\%\) pode ser construído por
\[
\widehat{P}
\pm
1.96\,\text{EP}.
\]

Portanto, para reduzir o erro pela metade, é necessário quadruplicar o número de simulações. Essa característica é uma propriedade fundamental do método de Monte Carlo.

\section{Aplicação ao Problema da Rede}

No código implementado, cada simulação é composta pelas seguintes etapas:

\begin{enumerate}
    \item Geração aleatória de um novo conjunto de resistências, simulando entupimentos;
    \item Construção das matrizes da rede;
    \item Resolução do sistema linear esparso;
    \item Cálculo da pressão máxima;
    \item Verificação se a pressão máxima excede o limite admissível.
\end{enumerate}

Esse procedimento é repetido \(N=1000\) vezes em cada bloco de simulação. Ao final de cada bloco, são calculadas estatísticas como:

\begin{itemize}
    \item pressão máxima observada;
    \item pressão média;
    \item probabilidade estimada de exceder o limite.
\end{itemize}

O processo completo é repetido em vários blocos independentes, permitindo avaliar a estabilidade das estimativas e quantificar a variabilidade entre execuções.

\section{Pseudocódigo do Algoritmo}

\begin{verbatim}
Entrada:
    Rede base
    Resistências R
    Probabilidade de entupimento p
    Fator de aumento alpha
    Limite de pressão L
    Número de simulações N

Para k = 1 até N:
    Gerar R^(k) aleatoriamente
    Resolver a rede com R^(k)
    Calcular Pmax^(k)

    Se Pmax^(k) > L:
        contador = contador + 1

Fim

Estimativa da pressão máxima média:
    media = média(Pmax^(1), ..., Pmax^(N))

Estimativa da probabilidade:
    prob = contador / N
\end{verbatim}

\section{Conclusão}

O método de Monte Carlo é adequado para analisar redes sujeitas a incertezas aleatórias, como entupimentos probabilísticos em condutos. Sua implementação permite estimar tanto o valor esperado da pressão máxima quanto a probabilidade de ocorrência de sobrepressões acima de um limite especificado. A precisão das estimativas aumenta com o número de simulações, com erro decrescendo proporcionalmente a \(1/\sqrt{N}\).

\end{document}